\documentclass[12pt]{article}
\usepackage{amsmath,amssymb,amsfonts}
\usepackage[margin=1in]{geometry}

\begin{document}

\textbf{Chapter 6, Problem 24.} \textit{An alternative approach to Hilbert transforms.}

(a) Show that if $f(z)$ is analytic in the upper half $z$-plane, then if $a = \alpha$ ($\alpha > 0$),
\[
f(a) = \frac{1}{\pi i}\int_{-\infty}^{\infty}\frac{f(x)}{x - a}\,dx,
\]
where $C$ is a semicircle of arbitrarily large radius with the real axis for its base.

(b) If $f(z) \to 0$ as $|z| \to \infty$ on the upper half-plane, then show that
\[
f(a) = \frac{1}{\pi i}\mathcal{P}\!\!\int_{-\infty}^{\infty}\frac{f(x)}{x-\alpha}\,dx \quad \text{for real } \alpha.
\]

(c) If $f(z) = u(x,y) + iv(x,y)$, show that (b) reduces to
\[
u(\alpha,0) = \frac{1}{\pi}\mathcal{P}\!\!\int_{-\infty}^{\infty}\frac{v(x,0)}{x-\alpha}\,dx, \qquad v(\alpha,0) = -\frac{1}{\pi}\mathcal{P}\!\!\int_{-\infty}^{\infty}\frac{u(x,0)}{x-\alpha}\,dx,
\]
which is the Hilbert transform pair.

(d) Hence show that if we let $\beta = 0$,
\[
u(\alpha,0) = \frac{1}{\pi}\mathcal{P}\!\!\int_{-\infty}^{\infty}\frac{v(x,0)}{x-\alpha}\,dx, \qquad v(\alpha,0) = -\frac{1}{\pi}\mathcal{P}\!\!\int_{-\infty}^{\infty}\frac{u(x,0)}{x-\alpha}\,dx.
\]

\bigskip
\textbf{Solution.}

\textbf{(a)} Consider the contour consisting of the real axis from $-R$ to $R$ and a semicircular arc $\Gamma_R$ in the upper half-plane. Since $f(z)$ is analytic in the upper half-plane and $a$ is in the upper half-plane ($\text{Im}(a) > 0$), by Cauchy's integral formula:
\[
f(a) = \frac{1}{2\pi i}\oint_C \frac{f(z)}{z-a}\,dz = \frac{1}{2\pi i}\left[\int_{-R}^{R}\frac{f(x)}{x-a}\,dx + \int_{\Gamma_R}\frac{f(z)}{z-a}\,dz\right].
\]

If $f(z) \to 0$ fast enough as $|z| \to \infty$ in the upper half-plane, then $\int_{\Gamma_R} \to 0$ as $R \to \infty$, giving:
\[
f(a) = \frac{1}{2\pi i}\int_{-\infty}^{\infty}\frac{f(x)}{x-a}\,dx.
\]

\textbf{(b)} Now let $a = \alpha$ be real. The integrand has a pole on the real axis. We use a small semicircular indentation of radius $\epsilon$ above the pole at $x = \alpha$. The contour becomes: real axis from $-R$ to $\alpha-\epsilon$, small semicircle $\gamma_\epsilon$ (going above the pole), real axis from $\alpha+\epsilon$ to $R$, and the large semicircle $\Gamma_R$.

Since $\alpha$ is now on the real axis (boundary), there's no pole inside the contour. But indenting into the upper half-plane, $\alpha$ is excluded. By Cauchy's theorem the integral is zero (no poles inside). The small semicircle contributes:
\[
\int_{\gamma_\epsilon}\frac{f(z)}{z-\alpha}\,dz \to -\pi i\,f(\alpha) \quad \text{as } \epsilon \to 0,
\]
(half-residue with negative sign since we go clockwise around half the pole).

So:
\[
0 = \mathcal{P}\!\!\int_{-\infty}^{\infty}\frac{f(x)}{x-\alpha}\,dx - \pi i\,f(\alpha) + 0.
\]
\[
f(\alpha) = \frac{1}{\pi i}\mathcal{P}\!\!\int_{-\infty}^{\infty}\frac{f(x)}{x-\alpha}\,dx. \quad \checkmark
\]

\textbf{(c)} Write $f = u + iv$ on the real axis. Then:
\[
u(\alpha) + iv(\alpha) = \frac{1}{\pi i}\mathcal{P}\!\!\int_{-\infty}^{\infty}\frac{u(x)+iv(x)}{x-\alpha}\,dx.
\]

Multiply $1/(\pi i) = -i/\pi$:
\[
u(\alpha) + iv(\alpha) = \frac{-i}{\pi}\mathcal{P}\!\!\int_{-\infty}^{\infty}\frac{u(x)+iv(x)}{x-\alpha}\,dx = \frac{1}{\pi}\mathcal{P}\!\!\int_{-\infty}^{\infty}\frac{v(x)}{x-\alpha}\,dx - \frac{i}{\pi}\mathcal{P}\!\!\int_{-\infty}^{\infty}\frac{u(x)}{x-\alpha}\,dx.
\]

Equating real and imaginary parts:
\[
\boxed{u(\alpha) = \frac{1}{\pi}\mathcal{P}\!\!\int_{-\infty}^{\infty}\frac{v(x)}{x-\alpha}\,dx, \qquad v(\alpha) = -\frac{1}{\pi}\mathcal{P}\!\!\int_{-\infty}^{\infty}\frac{u(x)}{x-\alpha}\,dx.}
\]

These are precisely the Hilbert transform relations.

\end{document}
